\chapter{Introduction}
\label{ch:introduction}

% \section{Problem Statement}


Emotions are a very important factor that contributes to video game enjoyment. 
For many game studios, such as \parencite{11BitStudios2026}, they are the core focus when designing games.
The people who create games make design decisions that shape the emotions that players feel.
Voice input could potentially be one of the things that can have an impact on the player's emotions. 
This thesis explores that possibility.
The goal of this paper is to verify how the player feels when they navigate in a sample video game environment using spoken natural language compared to traditional input methods, such as keyboard, mouse or controller.  
These methods rely on hands for input and control, making it impossible for handicapped people who lack the motor
skills and/or can't use their hands to take part in the activity. While games that support voice commands
do exist \parencite{iridiumstudiosThereCameEcho2015},
they rely on pre-existing voice commands. This poses a challenge by requiring the user to remember the exact commands
that need to be said to trigger actions. The system in this paper will utilise a language model to translate natural
language into specific game actions, thereby eliminating the need for command memorisation. 

Putting above in more technical terms with the current AI transcribing abilities it's possible to detect and transcribe
human voice in near real time \parencite{OpenAIWhisper2022}. That natural language transcription would be interpreted by an \textcite{LargeLanguageModel2025} with direct access to
the game engine. Such deep integration could produce an input method that’s more flexible than traditional voice commands
for navigation. 

This paper will focus on a game environment similar to a pen and paper \textcite{RPG2025} games, since players tend to have a variety
of options available to them at any time, promoting non-linear thinking and ambiguous commands.
This would allow for deep exploration into interpreting natural language as game commands. 

The technical challenge for this paper is: creating a natural voice input pipeline that can interpret and execute player
commands in the provided gaming environment with a high accuracy. This involves addressing software, usability,
and hardware challenges. 


\section{Research questions/objectives}

This thesis will follow two main research questions:

\begin{enumerate}
    \item How would the participants describe the experience of using voice input in the presented game environment? 
    \item How do users feel about using the voice modality compared to traditional input? 
\end{enumerate}


\section{Research Strategy}

% none of it is a reason for employing qualitative methods. qualitative methods are valid for collecting data on user's attitudes towards the assessed tool. You do not do it because there is no time, you do it because you want to understand the user's impressions and opinions on the presented intervention. Please rephrase and reference sources about the qualitative methods you deployed.
 
The research questions focus on emotions and personal opinions. For this reason the thesis utilises qualitative methods to seek and find the views of the participants. 
The methods used are going to be \parencite{youngImprovingLibraryUser2014}. Players will play both versions of the game, 10 minutes for each version.
After that, they will get interviewed with open-ended questions asking for their opinions and feelings towards both game versions. The interview part is structured after \parencite{mannResearchInterview2016}.


% Given the personal nature of this topic, and the limitations on time and resources, this thesis will be focusing
% on employing qualitative methods. This allows for empathising with the users on an individual level to understand
% the nuance of their experiences with the voice input method. 

% The main qualitative method used is going to be guided interviews with the subjects. That will seek out the 
% emotions and insights from the subjects. 

% This strategy, relying on qualitative will allow for reaching results that will be a combination of voices of
% all participants. 

% \section{Scope and limitations}


% "Chapter 2 introduces the theoretical foundations and relevant literature underpinning this work. The methodology behind the user testing phase is then detailed in Chapter 3, covering participant selection, session design and analytical approach. Chapter 4 documents the technical implementation..."

\section{Thesis structure overview}

% The structure for this thesis will follow the following chapters:

% As a first section, Chapter~\ref{ch:introduction}. It focuses on outlining the thesis and it's goals. 

% Continuing, Chapter~\ref{ch:review}. That chapter focuses on introducing the necessary definitions and explores the scientific literature relevant to this thesis.

% Following that comes Chapter~\ref{ch:method}. This part of the document is about the details of the user testing phase, explaining the full process of gathering user responses.

% Right after that is the Chapter~\ref{ch:implementation}. That section is dedicated to explaining the technical details of the project. It goes through every step that allows the game to work.

% Next in order is Chapter~\ref{ch:results}. It outlines the five main findings from the user testing with supporting quotes for each of them.

% Building on those findings comes Chapter~\ref{ch:discussion}. It seeks to explain what the findings mean and how they fit with the existing literature. 

% Lastly comes \ref{ch:conclusion}. It gathers all of the observations and suggestions and directly tackles the research questions.

% In addition comes the \ref{appendix_a}, in which the Design Document of the created video game is laid out. 

Chapter~\ref{ch:introduction} introduces the thesis and showcases the scope of the work. 
Following that, the Chapter~\ref{ch:review} provides the necessary definitions and examines the scientific literature relevant to this thesis.
The methodology of the user testing is described in detail in Chapter~\ref{ch:method}, covering participant information, user-testing details and the full process of gathering responses.
Right after that is the Chapter~\ref{ch:implementation}. That part of the thesis focuses on the technical details of the game prototype, diving deep into the inner workings of the game and the design process. 
Chapter~\ref{ch:results} documents the results from the user-testing aggregated into five main categories.
Following that, Chapter~\ref{ch:discussion} takes those five categories and explains their meaning, analysing both the potential implications and the connection to existing literature.
\ref{ch:conclusion} finishes that process. It takes all of the observations and uses them to answer the research questions presented in this thesis. 
As an addition, Appendix~\ref{appendix_a} includes the design document of the video game made as a part of this thesis.






% \begin{enumerate}
%     \item Introduction
%     \item Literature Review
%     \item Methodology
%     \item Implementation
%     \item Results
%     \item Discussion
%     \item Conclusion
%     \item Appendix
% \end{enumerate}